% Original introduction (preserved for reference):
% Electricity demand has been increasing at unprecedented rates, with projections indicating a compound annual rate of 3.6\% over 2026-2030 -- approximately 50\%  than that of the preceding decade \cite{iea_electricity}.
% A significant contributor to this demand is the rapid expansion of datacenters, which current currently consume 1.5\% of global electricity and are projected to increase by 15\% annually through 2030, a rate four times greater than that of any other sector \cite{iea_demand_report}.
% A significant driver for this growth is artificial intelligence (AI), with model training and inference taking up majority of the compute jobs.
% Previous efforts for reducing carbon emissions centered on improving energy efficiency, measured through Power Usage Effectiveness (PUE).
% By today's standard, most large datacenters are close to 1.0 PUE with Google reporting their large-scale data centers with a PUE of 1.09, suggesting server-level energy efficiency is approaching practical limits \cite{google_power_usage_effectiveness}.
%
% Electric vehicles (EVs) are achieving widespread adoption across households globally and are a growing workload on electricity grids.
% EV charging consumed 180 TWhs in 2024, and is projected to triple by 2030, representing 2.5\% of the global electricity demand \cite{iea_ev_report}.
% Given these trends, the imperative to achieve carbon efficiency has become increasingly urgent.
% It is important to distinguish energy efficiency is not equivalent to carbon efficiency.
% A highly energy-efficient workload operating on a carbon-intensive grid will not achieve carbon efficiency.
% Conversely, a workload with poor energy efficiency running on a fully renewable grid is inherently carbon efficient, irrespective of its absolute energy consumption.
%
% An established mechanism for pursuing carbon efficiency is the use of carbon offset accounting, where certificates are purchased to fund the development clean energy infrastructure -- such as solar and wind projects -- and in return the purchaser will receive carbon offsets \cite{maji_green_2024}.
% While this framework has supported electricity grids transition toward being more carbon efficient, it brings the challenge of managing the inherent intermittency of renewable energy.
% Renewable energy, such as wind and solar, are highly dependant on the weather and thus controlling costs and planning to utilize this renewable energy requires high effort.
%
% To address this challenge, researchers have increasingly explored strategies to improve the carbon efficiency of workloads directly.
% To reach a carbon-efficient outcome, coordinated efforts on both the supply and demand side of the electric system is required.
% On the supply side, organizations can enter into power purchase agreements (PPAs) for renewable energy to power large-scale operations such as data centers \cite{maji_green_2024}.
% On the demand side, techniques are being developed to modulate a workload's electricity consumption -- and consequently its carbon footprint -- by shifting execution to times and/or locations.
% For instance, EV charging sessions can be rescheduled to coincide with periods of minimal grid carbon intensity, a strategy known as temporal shifting.
% Cloud-based compute jobs are a workload that can exploit "spatial shifting", as it can be migrated across a geographically distributed cluster to access grids with the lowest carbon intensity.
% The nature of the workload itself further influences the available optimization strategies.
% A continuous workload, such as a stateful compute job, will execute for the duration of the job and cannot be shifted of interrupted during execution.
% An interruptable workload, such as EV charging, can be paused and resumed during execution time without affecting the final outcome.
% Temporal and spatial shifting, combined with the distinction between continuous and interruptible workloads, constitute the principal degrees of freedom that carbon-aware optimizers exploit to minimize operational carbon emissions \cite{aslan_toward_2024,google_data_center_work_harder,gsteiger_caribou_2024,souza_casper_2023,sukprasert_limitations_2024}.
%
% To achieve carbon optimality, workloads use a carbon-aware optimizer to make the temporal/spatial shifts to achieve carbon optimality.
% These carbon-aware optimizers depend on carbon forecasts of grids to inform their decision-making \cite{grid_observability}.
% Current state-of-the-art forecasting approaches construct data-intensive and computationally expensive neural networks to capture the complex dynamics of a grid's energy source mix  \cite{maji_carboncast_2022,yan_ensembleci_2025,zhang_gnn-based_2023}.
% Although these models can achieve high accuracy in their forecasts, their effectiveness in scheduling workloads to be carbon optimal is not explored.
% Moreover, the practical applicability of these models is constrained by their reliance on extensive historical training data and costly hardware requirements, rendering real-time, adaptive decision-making more complex and time-consuming in practice.
% With electricity grids constantly evolving and changing their energy sources mix, it is critical to capture these changes when forecasting to make the best decision.
% In this thesis, we propose utilizing an alternative evaluation metric -- concordance -- to assess the effectiveness of carbon intensity forecasts in the context of workload scheduling.
% Concordance measures the degree to which the ordinal ranking of forecasted values preserves that of the realized carbon intensity observations.
% We integrate this metric into our own carbon forecasting model, LiteCast.
% LiteCast is a lightweight time-series model that aims for high concordance and optimal carbon scheduling.

Electricity demand has been increasing at unprecedented rates, with projections indicating a compound annual growth rate of 3.6\% over 2026--2030, which is approximately 50\% higher than that of the preceding decade \cite{iea_electricity}.
A significant contributor to this demand is the rapid expansion of data centers, which currently consume 1.5\% of global electricity and are projected to grow by 15\% annually through 2030, a rate four times greater than that of any other sector \cite{iea_demand_report}.
A primary driver of this growth is artificial intelligence (AI), with model training and inference consuming the majority of data center compute.
Electric vehicles (EVs) are also a rapidly growing source of electricity demand, with EV charging consuming 180~TWh in 2024 and projected to triple by 2030, representing 2.5\% of global electricity demand \cite{iea_ev_report}.
Together, these trends are placing unprecedented pressure on electricity grids, intensifying the need to address their environmental impact.

Previous efforts to reduce the environmental impact of computing have centered on improving energy efficiency, commonly measured through Power Usage Effectiveness (PUE).
PUE measures the ratio of total facility energy to IT equipment energy, where a value of 1.0 represents perfect efficiency.
By today's standards, most large data centers operate close to this ideal, with Google reporting a PUE of 1.09 for its large-scale facilities \cite{google_power_usage_effectiveness}, suggesting that server-level energy efficiency is approaching its practical limits.
Critically, high energy efficiency is not equal to high carbon optimality.
A highly energy-efficient workload operating on a carbon-intensive grid will not achieve carbon optimality, while a workload with poor energy efficiency running on a fully renewable grid is inherently carbon optimal.
Given the diminishing returns of energy efficiency improvements, the imperative to pursue carbon optimality has become increasingly urgent.

Electricity grids draw from a mix of energy sources, including fossil fuels, renewables, and low-carbon sources, each with different carbon factors.
The carbon intensity of the grid fluctuates over time and across regions depending on which sources are generating power at any given moment.
Carbon-aware optimizers exploit this variability by shifting workloads to times and locations where the carbon intensity is lowest, reducing their operational carbon footprint \cite{google_data_center_work_harder,gsteiger_caribou_2024,souza_casper_2023,sukprasert_limitations_2024}.
Two principal strategies exist: \textit{temporal shifting}, which defers or interrupts a workload to coincide with periods of lower carbon intensity, and \textit{spatial shifting}, which migrates a workload to a region with a cleaner grid.
The applicability of these strategies depends on the nature of the workload: a continuous workload must execute uninterrupted once started, while an interruptible workload can be paused and resumed without affecting its outcome.
For example, a batch compute job is a continuous workload, while EV charging is an interruptible one.

Carbon-aware optimization requires planning ahead, as a scheduler must know when and where the carbon intensity will be lowest before a workload begins.
This planning depends on carbon intensity forecasts, making forecasting a critical enabler of carbon-aware optimization \cite{grid_observability}.
Previous forecasting efforts have focused on maximizing accuracy by minimizing error metrics such as mean absolute percentage error (MAPE), building data-intensive and computationally expensive neural networks \cite{maji_carboncast_2022,yan_ensembleci_2025,zhang_gnn-based_2023}.
However, it has not been examined whether higher forecast accuracy actually produces better scheduling outcomes.

In this thesis, we show that what drives better scheduling decisions is not only accuracy, but \textit{concordance}: the degree to which a forecast preserves the correct ordering of carbon intensity values over time.
A scheduler that correctly identifies the lowest-carbon window will achieve better carbon outcomes, even if the predicted values are not precise.
To demonstrate this, we present \textit{LiteCast}, a lightweight carbon intensity forecasting model that prioritizes concordance, requires less historical data than state-of-the-art models, runs on basic compute hardware, and adapts quickly to changes in the electricity grid.
We evaluate LiteCast across 48 regions worldwide under a range of workloads (including synthetic batch jobs and a real supercomputer cluster trace), demonstrating near-optimal carbon scheduling performance.
